\documentclass[a4paper]{arrowhead}

\usepackage[yyyymmdd]{datetime}
\usepackage{etoolbox}
\usepackage[utf8]{inputenc}
\usepackage[T1]{fontenc}
\usepackage{multirow}
\usepackage{longtable}
\usepackage{tikz}
\usetikzlibrary{positioning,shapes.multipart,arrows.meta}

\renewcommand{\dateseparator}{-}
\setlength{\emergencystretch}{3em}
\lstset{language=HTTP,breakatwhitespace=false,columns=fullflexible}

%% Special references
\newcommand{\mref}[2][]{{\textcolor{ArrowheadPurple}{\hyperref[sec:model:#2]{\ifstrempty{#1}{#2}{#1}}}}}
\newcommand{\oref}[1]{{\textcolor{ArrowheadBlue}{\hyperref[sec:operations:#1]{#1}}}}
\newcommand{\pdef}[1]{{\textcolor{ArrowheadGrey}{#1\label{sec:model:primitives:#1}\label{sec:model:primitives:#1s}\label{sec:model:primitives:#1es}}}}
\newcommand{\pref}[1]{{\textcolor{ArrowheadGrey}{\hyperref[sec:model:primitives:#1]{#1}}}}

% as template_sd.tex; the input and output arguments are passed already formatted (\mref or \pref)
\newrobustcmd\osubsection[3]{
  \addtocounter{subsection}{1}
  \addcontentsline{toc}{subsection}{\protect\numberline{\thesubsection}operation \textcolor{ArrowheadBlue}{#1}}
  \subsection*{
    \thesubsection\quad
    operation
    \textcolor{ArrowheadBlue}{#1}
    (\notblank{#2}{#2}{})
    \notblank{#3}{: #3}{}
  }
  \label{sec:operations:#1}
}
\newrobustcmd\msubsection[2]{
  \addtocounter{subsection}{1}
  \addcontentsline{toc}{subsection}{\protect\numberline{\thesubsection}#1 \textcolor{ArrowheadPurple}{#2}}
  \subsection*{\thesubsection\quad#1 \textcolor{ArrowheadPurple}{#2}}
  \label{sec:model:#2} \label{sec:model:#2s} \label{sec:model:#2es}
}
%%

\usepackage{xurl}
\begin{document}
\sloppy

%% Arrowhead Document Properties
\ArrowheadTitle{sysMonitor}
\ArrowheadServiceID{sysMonitor}
\ArrowheadType{Service Description}
\ArrowheadTypeShort{SD}
\ArrowheadVersion{5.3.0}
\ArrowheadDate{2026-11-09}
\ArrowheadAuthor{Jerker Delsing}
\ArrowheadStatus{REVIEW DRAFT}
\ArrowheadContact{jerker.delsing@ltu.se}
\ArrowheadFooter{\href{www.arrowhead.eu}{www.arrowhead.eu}}
\ArrowheadSetup
%%

%% Front Page
\begin{center}
  \vspace*{1cm}
  \huge{\arrowtitle}

  \vspace*{0.2cm}
  \LARGE{\arrowtype}
  \vspace*{1cm}

  \Large{Service ID: \textit{"\arrowid"}}

  \vspace*{\fill}

  % Front Page Image

  \vspace*{1cm}
  \vspace*{\fill}

  % Front Page Abstract
  \begin{abstract}
    This document describes the native \texttt{sysMonitor} service, version 1.0.0, produced by the AIGym support system of Arrowhead (documentation edition 5.3.0).
  \end{abstract}

  \vspace*{1cm}
\end{center}

\newpage
%%

%% Table of Contents
\tableofcontents
\newpage
%%

\section{Documentation edition and implementation scope}
This is the 5.3.0 documentation edition, dated 2026-11-09. It carries forward the system and service contracts from the retained 5.2.2 documentation and identifies the current primary implementation sources. The former \path|documentation/v5.2.2| folder is retained unchanged.

Document version 5.3.0 is distinct from executable, dependency, registered system/service and wire-contract versions. Values still defined by the source (including 5.2.2 artefact versions, 1.0.0 service versions and \path|SysMonitor-AMQP-IDD-5.2.2| payload identifiers) are preserved. This documentation update does not change runtime behaviour or establish a new implementation release. Historical verification statements describe the tests recorded in the inherited baseline; this revision validates document compilation and references.

\section{Overview}
\label{sec:overview}

This document describes the \texttt{sysMonitor} service, which provides native operational telemetry of the AIGym system: periodic \texttt{echo}, \texttt{power} and \texttt{status} observations published to the local cloud's AMQP broker.
The service is produced by AIGym \cite{aig-sysd}; its concrete interface is the AMQP IDD \cite{aig-sysMonitor-amqp-idd}.

The rest of this document is organized as follows.
In Section \ref{sec:operations}, we describe the abstract message operations provided by the service.
In Section \ref{sec:model}, we end the document by presenting the data types used by the mentioned operations.

\subsection{Significant Prior Art}

Arrowhead 4.x and upstream 5.2.2 had an HTTP \texttt{monitor}/\texttt{echo} endpoint. In this workspace the telemetry follows SysMonitor-AMQP-IDD-5.2.2; the HTTP liveness endpoint is not registered.

\subsection{How This Service Is Meant to Be Used}

The SysMonitor support system subscribes to the observations (selected with \texttt{sys\allowbreak{}Monitor\allowbreak{}Management}) and keeps the history; LocalCloudManager shows it.

\subsection{Important Delimitations and Dependencies}
\label{sec:delimitations}

The observations describe the process, not the correctness of its answers. The service is offered only with \texttt{sysmonitor.\allowbreak{}enabled=\allowbreak{}true}. Delivery is publisher-confirmed but not end-to-end acknowledged by SysMonitor.

\newpage

\section{Service Interface}
\label{sec:operations}

The \texttt{sysMonitor} service is a publishing service: the producer publishes one observation per operation (observation type) every publication interval (default 30 s) and on every event of that type; consumers receive them through the broker. Each operation has no input and the output \mref{Observation}.

\begin{figure}[ht!]
  \centering
  \begin{tikzpicture}[font=\small]
    \node[draw=ArrowheadBlue, very thick, rectangle split, rectangle split parts=2, align=center,
          rectangle split part fill={ArrowheadBlue!12,white}, text width=9cm, inner sep=6pt] (b)
      {\guillemotleft{}Service\guillemotright{}\\\textbf{sysMonitor}
        \nodepart{second} \texttt{echo}\\ \texttt{power}\\ \texttt{status}\\ \texttt{trainingSample}\\ \texttt{inference}};
  \end{tikzpicture}
  \caption{SysML-style block of the \texttt{sysMonitor} service and its observation types.}
  \label{fig:SD}
\end{figure}

\osubsection{echo}{}{\mref{Observation}}

Liveness: \texttt{status} (\texttt{alive} or \texttt{degraded}), \texttt{uptimeSeconds} and a \texttt{detail} text; quality \texttt{good} or \texttt{uncertain}.

\osubsection{power}{}{\mref{Observation}}

Power supply: \texttt{supplyType=\allowbreak{}line}, \texttt{state=\allowbreak{}available}; unit \texttt{1}.

\osubsection{status}{}{\mref{Observation}}

Process state: \texttt{state=\allowbreak{}running}, \texttt{health} (\texttt{healthy}/\texttt{degraded}) and \texttt{serviceCount}, the number of registered services.

\osubsection{trainingSample}{}{\mref{Observation}}

Published for every added sample: the sample's target value, with model and consumer in the metadata.

\osubsection{inference}{}{\mref{Observation}}

Published for every prediction: the predicted value, with model, consumer and neighbours in the metadata.

\newpage

\section{Information Model}
\label{sec:model}

\msubsection{struct}{Observation}

One observation.

\noindent\begin{tabularx}{\textwidth}{| p{4.25cm} | p{3.5cm} | X |} \hline
\rowcolor{gray!33} Field & Type & Description \\ \hline
schemaVersion & \pref{Version} & Version of the observation schema, \texttt{1.0}. \\ \hline
observationId & \pref{Identifier} & Unique identifier of the observation. \\ \hline
providerSystem & \pref{Name} & The publishing system. \\ \hline
serviceInstance & \pref{Name} & The service instance (default the system name). \\ \hline
observationType & \pref{Name} & \texttt{echo}, \texttt{power}, \texttt{status}, \texttt{trainingSample}, \texttt{inference}. \\ \hline
sourceTime & \pref{DateTime} & Time of the observation. \\ \hline
sequence & \pref{Number} & Sequence number per publisher. \\ \hline
value & \pref{Metadata} & The observed value (per operation). \\ \hline
unit & \pref{Name} & Unit of the value, if any. \\ \hline
quality & \pref{Name} & \texttt{good} or \texttt{uncertain}. \\ \hline
metadata & \pref{Metadata} & \texttt{publisherMode=native}. \\ \hline
\end{tabularx}

\subsection{Primitives}
\label{sec:model:primitives}

\noindent\begin{tabularx}{\textwidth}{| p{3cm} | X |} \hline
\rowcolor{gray!33} Type & Description \\ \hline
\pdef{DateTime} & Pinpoints a specific moment in time (UTC). \\ \hline
\pdef{Identifier} & A string that identifies a record, e.g. a UUID or a composite identifier. \\ \hline
\pdef{Metadata} & A set of key-value pairs with arbitrary, possibly nested values. \\ \hline
\pdef{Name} & A string identifier that is intended to be both human and machine-readable. \\ \hline
\pdef{Number} & An integer number. \\ \hline
\pdef{Version} & Specifies a version, semantic versioning \texttt{major.minor.patch}. \\ \hline
\end{tabularx}

\newpage

\bibliographystyle{IEEEtran}
\bibliography{bibliography}

\newpage

\section{Revision History}
\subsection{Amendments}

\noindent\begin{tabularx}{\textwidth}{| p{1cm} | p{3cm} | p{2cm} | X | p{4cm} |} \hline
\rowcolor{gray!33} No. & Date & Version & Subject of Amendments & Author \\ \hline

1 & 2026-10-22 & 5.2.2 & Document created from the Arrowhead LaTeX templates for the 5.2.2 Java Spring implementation, including the local extensions of this workspace. & Jerker Delsing \\ \hline
2 & 2026-10-30 & 5.2.2 & Checked against the primary Java source tree (Eclipse-Arrowhead-java): services, operations, paths and data models unchanged. & Claude Opus 5.5 (draft) \\ \hline
3 & 2026-11-09 & 5.3.0 & New documentation edition; primary-source references and Go/Rust implementation scope; 5.2.2 retained. & Documentation maintenance \\ \hline
\end{tabularx}

\subsection{Quality Assurance}

\noindent\begin{tabularx}{\textwidth}{| p{1cm} | p{3cm} | p{2cm} | X |} \hline
\rowcolor{gray!33} No. & Date & Version & Approved by \\ \hline

1 & & \arrowversion  &  \\ \hline

\end{tabularx}

\end{document}
